\textbf{Catálogo y demanda de integración.}
La oferta selecciona diez interfaces funcionales externas: contabilidad, acceso, CCTV, meteorología, surtidor, SMS, mensajería instantánea, correo, navegación de baja capacidad y extracción documental. Los servicios de plataforma AWS, UEM y observabilidad conservan su inventario propio en T-11; no se cuentan nuevamente como interfaces de negocio. MQTT de medidores pertenece a la captura interna de telemetría. DIRECTEMAR y VHF conservan sus canales de autoridad y coordinación humana: no se les atribuye una API ni se cuentan como una undécima integración automática. El catálogo A4.2 de \ReferenciaAnexo{Anexos de SD4}{SYNAPTIX-Subdocumento4-Anexos.pdf} fija flujo, contrato, autenticación, volumen, contingencia y prueba por interfaz.

En un mes de treinta días, los supuestos del catálogo producen 402.120 mensajes lógicos: solicitudes, respuestas, cambios de estado y acuses de negocio, sin contar segmentos TCP ni el video ordinario. Para contabilidad se usan 2.300 documentos/mes y 296 consultas diarias de estado; acceso usa 4.000 personas/mes con entrada y salida; meteorología una observación/minuto. Los demás coeficientes son hipótesis explícitas de utilización de Synaptix. Las tasas de peak se prueban por separado del promedio mensual y se concilian con EST-11 de SD2 y SD5, 5.4.1. Una API de proveedor se selecciona mediante su documentación, mientras la compatibilidad del sistema existente se verifica sobre su variante real.

Para los tres canales de emergencia se seleccionan Twilio Programmable Messaging para SMS, Meta WhatsApp Cloud API para mensajería y Amazon SES API v2 para correo. Cada canal reserva cinco envíos/s; con un reintento por destinatario, una campaña a 296 personas requiere $296\times3\times2=1.776$ intentos, 592 por canal y $592/5=118,4$ segundos de servicio por canal, antes de pausas y errores. La aceptación conserva el límite de diez minutos hasta el último envío, ráfagas, fallas, acuses y escalamiento a llamada registrada. Entrega técnica no equivale a acuse humano. El enrolamiento verifica los tres destinos, autorización, plantillas y cuotas contratadas antes de habilitar; la selección del proveedor no representa un contrato suscrito ni elimina la brecha de una persona sin canal habilitado \parencites{c8-twilio-mensajes,c8-meta-mensajes,c8-ses-eventos}[cap.~15, RT-16.21]{caso07}.

\textbf{Navegación de baja capacidad.}
La aplicación Kotlin selecciona HTTPS con una cola durable y un endpoint acotado de intercambio, utilizable por datos móviles intermitentes o por IP satelital. Iridium GO! exec con Certus 100 es la referencia satelital de interoperabilidad: ofrece acceso IP mediante Wi-Fi; no se utiliza como SMS ni se presupone que el equipo ya llevado por cada armador sea esa variante. El perfil del equipo habilita únicamente DNS y HTTPS necesarios y bloquea tráfico ajeno al intercambio. La compatibilidad, suscripción y pruebas se registran por equipo voluntariamente enrolado; no se instala seguimiento automático a bordo. Una conexión incompatible no se presenta como canal habilitado \parencites{c8-iridium-exec,c8-iridium-faq}.

El contrato fija 1 KiB máximo por mensaje de negocio UTF-8, incluida su envolvente; el intercambio completo reserva hasta 16 KiB de transporte, autenticación y respuesta. Se usan identificador de plan y evento, secuencia, origen, tiempo de actuación, vencimiento, tipo, datos mínimos y huella; los detalles personales y anexos se cargan antes de zarpar. Una ruta extensa se divide en mensajes referenciados, cada uno dentro del límite, con confirmación del manifiesto. No se incluyen fotografías, salud de tripulantes ni un envío periódico de posición por defecto. Se reserva una hipótesis de ocho mensajes por sentido y embarcación/día, incluidos acuses y avisos: para 296 embarcaciones, $296\times16=4.736$ mensajes/día. El intercambio programado ocurre cada seis horas mientras se utiliza ese modo; declarar, corregir o cerrar un plan provoca un intento inmediato y no espera esa ventana. El usuario abre la sesión satelital cuando corresponde; la cola no presume conectividad automática del terminal.

La demora operacional tolerada para mensajes ordinarios es de hasta seis horas desde su generación; a las seis horas sin acuse se informa demora y se escala a Torre, y a las doce se marca no entregado y se conserva para conciliación. Una instrucción urgente no utiliza ese plazo como garantía de seguridad: continúa por VHF y los procedimientos de autoridad. El vencimiento del plan activa el escalamiento del caso en quince minutos aunque falte un mensaje de cierre. Los reintentos conservan identificador, con espera de 1, 5, 15 y 60 minutos y luego cada hora mientras el mensaje sea vigente; no se descarta un hecho registrado al vencer su instrucción. El receptor deduplica, espera secuencias anteriores y distingue tiempo de actuación de recepción; un cierre tardío no elimina una alerta ya registrada. Se autentica la sesión al transmitir y se conserva la evidencia del consentimiento y autoría de la actuación offline; no se prolongan tokens vencidos para conseguir el envío. La prueba demora mensajes seis y doce horas, altera el orden, duplica y pierde respuestas, corta IP y revoca consentimiento. La latencia de horas es una tolerancia de diseño ante ventanas de enlace, no una prestación de latencia publicada por Iridium \parencite[cap.~14, sec.~14.2; cap.~17, sec.~17.2]{caso07}.
